
iv)
Supposons enfin @@M0@@ régulière et @@M1@@ lisse. Pour tout @@M2@@, on peut trouver un voisinage ouvert @@M3@@ de @@M4@@, un voisinage ouvert @@M5@@ de @@M6@@ au-dessus de @@M7@@, et un diagramme commutatif :
@@M8@@
où @@M9@@ est le morphisme structural, et @@M10@@ un morphisme étale (EGA IV 17.11.4). On peut alors supposer @@M11@@, @@M12@@. Il est facile de voir que @@M13@@ est une immersion (introduire le produit fibré @@M14@@). Appliquant les résultats ii) et iii), on en déduit que @@M15@@ est régulière, puis que @@M16@@ est régulière.

Montrons maintenant 1.2.

La propriété étant locale sur @@M17@@, on peut, pour la prouver, se placer au voisinage d'un point @@M18@@ de @@M19@@. On peut alors trouver un voisinage ouvert @@M20@@ de @@M21@@ dans @@M22@@, un @@M23@@-schéma lisse @@M24@@, et une @@M25@@-immersion régulière @@M26@@. On a alors le diagramme commutatif suivant:
@@M27@@

On en déduit une immersion fermée @@M28@@, et le diagramme commutatif :
@@M29@@
